\section{Evolución del paradigma de la Scaling Law}

La sección anterior trató el impacto de GPT y ChatGPT; esta entra en un paradigma técnico más fundamental. Al principio, la Scaling Law parecía un problema de escala en un solo eje; con el tiempo se convirtió en un problema de crecimiento de capacidades en varios ejes.

En la entrevista, «la evolución del paradigma de la Scaling Law» es uno de los ejes técnicos centrales. Zhang Peng (张鹏) señaló que el futuro no consiste solo en que siga mejorando la capacidad de los modelos fundacionales, sino también en la fusión de múltiples modelos y múltiples datos, los agentes y nuevas formas de cómputo; en particular, la línea de RL seguirá produciendo métodos nuevos. Es decir, la Scaling Law ya no equivale únicamente a la fórmula de preentrenamiento de «más parámetros, más datos, más cómputo».

\begin{figure}[H]
\centering
\includegraphics[width=0.96\textwidth]{scaling-law-evolution.png}
\caption{Evolución del paradigma de la Scaling Law: escala, multimodalidad, posentrenamiento, RL y Agent.}
\end{figure}

\begin{knowledgebox}{Lectura de la figura: la Scaling Law pasa de un eje a varios ejes}
El relato inicial de los grandes modelos ponía el acento en los parámetros, los datos y el cómputo; después se sumaron los datos multimodales, el posentrenamiento, la retroalimentación de preferencias y de tareas, el RL y los entornos de Agent. El crecimiento de capacidades no proviene solo de «ser más grande», sino también de «aprovechar mejor la retroalimentación».
\end{knowledgebox}

\subsection{Por qué el Agent se convirtió en la vía de aplicación}

Zhang Peng considera que los agentes son importantes porque resuelven el problema de cómo llevar un modelo hasta su aplicación real. El modelo es la base de capacidades, pero lo que el cliente quiere son resultados de negocio. El Agent conecta el modelo con herramientas, procesos, entornos y necesidades de los usuarios; es la capa intermedia para la comercialización y para liberar las capacidades.

\begin{knowledgebox}{Términos clave: Scaling Law y conceptos relacionados}
\begin{tabularx}{\textwidth}{p{0.23\textwidth}p{0.32\textwidth}X}
\toprule
Término & Explicación en una frase & Papel en este episodio \\
\midrule
Scaling Law & Regularidad con la que cambia la capacidad de un modelo al escalar parámetros, datos y cómputo & Relato de fondo del paradigma de los grandes modelos. \\
Post-training & Optimización continua del modelo después del preentrenamiento mediante preferencias, tareas y retroalimentación & Hace que el modelo sea más útil y más controlable. \\
RL & Reinforcement Learning, aprendizaje por refuerzo & Zhang Peng considera que traerá nuevas formas de cómputo. \\
Agent & Agente que invoca herramientas y ejecuta tareas & Vía de aplicación que conecta el modelo con usos reales. \\
Fusión de múltiples modelos & Varios modelos y fuentes de datos que en conjunto forman un sistema de capacidades & Una de las direcciones futuras de los modelos base. \\
\bottomrule
\end{tabularx}
\end{knowledgebox}

\subsection{Resumen de la sección}

La evolución de la Scaling Law muestra que la competencia entre grandes modelos ya no se reduce a «quién es más grande». La retroalimentación, el posentrenamiento, el Agent, el RL y los datos multimodales se están convirtiendo en nuevos ejes de crecimiento de capacidades.
